\section{Introduction}
\label{sec:intro}

When language-model text resembles human writing, its origin can be difficult to establish in
settings such as disinformation analysis, phishing investigation, and academic-integrity review
~\citep{liang2026survey}. A post-hoc detector learns recurring output patterns; a watermark instead
changes generation under a secret key and lets the key holder test a later document
~\citep{petitcolas1999hiding,kirchenbauer2023watermark}. We focus on watermarking because it can
offer a stated finite-sample false-positive rate, rather than relying on a classifier whose behaviour
may change with the decoder~\citep{krishna2023paraphrasing,
jovanovic2024stealing,zhao2024provable,christ2024undetectable}.

% Keep the first page available for prose, but prevent this top float from
% jumping above the title block on that page.
\suppressfloats[t]
% Conceptual contrast: the two panels differ in where the watermark signal lives.
% Concrete text replaces generic token strips so the distinction reads at a glance.
\begin{figure}[t]
\centering
\resizebox{\linewidth}{!}{%
\begin{tikzpicture}[
  x=1cm, y=1.08cm,
  font=\scriptsize,
  >=Stealth,
  paneltitle/.style={font=\scriptsize\bfseries, anchor=west, text=black!82},
  note/.style={font=\tiny, text=black!61, align=center},
  word/.style={draw=black!38, fill=white, rounded corners=0.7pt,
               minimum height=3.5mm, inner xsep=2.1pt, inner ysep=0.6pt,
               font=\tiny},
  keyed/.style={word, fill=black!20, draw=black!50},
  carrier/.style={word, fill=oursemph!7, draw=oursemph,
                  line width=0.75pt, text=oursemph},
  detector/.style={draw=black!38, fill=white, rounded corners=1.3pt,
                   minimum height=5.7mm, inner xsep=4pt,
                   font=\scriptsize\bfseries},
  bluedetector/.style={detector, draw=oursemph, fill=oursemph!6,
                       text=oursemph, line width=0.65pt},
  flow/.style={-{Stealth[length=4pt,width=4pt]}, draw=black!53,
               line width=0.5pt},
  blueflow/.style={-{Stealth[length=4pt,width=4pt]}, draw=oursemph,
                   line width=0.8pt},
  genflow/.style={-{Stealth[length=3pt,width=3pt]}, draw=black!53,
                  line width=0.45pt, shorten <=1pt, shorten >=1.2pt},
  bluegenflow/.style={-{Stealth[length=3pt,width=3pt]}, draw=oursemph,
                      line width=0.65pt, shorten <=1pt, shorten >=1.2pt},
  keyflow/.style={-{Stealth[length=3pt,width=3pt]}, dashed, line width=0.5pt,
                  shorten >=0.8pt}
]

\path[use as bounding box] (0,0) rectangle (13.88,-3.30);
\fill[black!2, rounded corners=2.3pt] (0.02,-0.02) rectangle (6.70,-3.18);
\fill[oursemph!4, rounded corners=2.3pt] (7.04,-0.02) rectangle (13.86,-3.18);

\node[paneltitle] at (0.22,-0.27) {(a) Signal in tokens};
\node[circle, draw=black!40, fill=white, minimum size=6.4mm,
      inner sep=0pt] (tokmodel) at (0.92,-1.15) {$\theta$};
\node[note] (tokkey) at (1.46,-0.67) {$K$ biases draws};
\node[word]  (t1) at (1.84,-1.15) {the};
\node[keyed] (t2) at (2.45,-1.15) {model};
\node[word]  (t3) at (3.20,-1.15) {left};
\node[word]  (t4) at (3.71,-1.15) {a};
\node[word]  (t5) at (4.22,-1.15) {clear};
\node[keyed] (t6) at (4.91,-1.15) {trace};
\draw[genflow] (tokmodel.east) -- (t1.west);
\draw[keyflow, draw=black!55] (tokkey.south) -- (1.46,-1.15);
\node[detector] (tokdet) at (3.37,-2.25) {keyed-token count};
\node[note, anchor=west] at (3.63,-1.72) {read $y,K$};
\draw[flow] (3.37,-1.54) -- (tokdet.north);
\node[detector] (tokp) at (5.55,-2.25) {$p$-value};
\node[note] at (5.55,-2.72) {no model call};
\draw[flow] (tokdet.east) -- (tokp.west);
\draw[decorate, decoration={brace, amplitude=3pt, mirror}, black!45]
  (1.55,-1.43) -- (5.20,-1.43);

\node[paneltitle, text=oursemph] at (7.24,-0.27) {(b) Signal in model response};
\node[circle, draw=oursemph, fill=oursemph!6, minimum size=6.4mm,
      inner sep=0pt, text=oursemph, line width=0.65pt]
  (ourmodel) at (7.92,-1.15) {$\theta$};
\node[note, text=oursemph] (ourkey) at (8.46,-0.67) {$K$ selects a response};
\node[word]    (s1) at (8.84,-1.15) {the};
\node[carrier] (s2) at (9.45,-1.15) {model};
\node[word]    (s3) at (10.20,-1.15) {left};
\node[word]    (s4) at (10.71,-1.15) {a};
\node[word]    (s5) at (11.22,-1.15) {clear};
\node[carrier] (s6) at (11.91,-1.15) {trace};
\draw[bluegenflow] (ourmodel.east) -- (s1.west);
\draw[keyflow, draw=oursemph] (ourkey.south) -- (8.46,-1.15);
\node[bluedetector] (ourdet) at (10.37,-2.25) {re-mask $+$ probe};
\node[note, text=oursemph, anchor=west] at (10.63,-1.72) {response signs};
\draw[blueflow] (10.37,-1.54) -- (ourdet.north);
\draw[blueflow] (ourmodel.south) to[out=-75,in=180] (ourdet.west);
\node[bluedetector] (ourp) at (12.62,-2.25) {$p$-value};
\node[note] at (12.62,-2.72) {rerun the same $\theta$};
\draw[blueflow] (ourdet.east) -- (ourp.west);
\draw[decorate, decoration={brace, amplitude=3pt, mirror}, oursemph]
  (8.55,-1.43) -- (12.20,-1.43);

\end{tikzpicture}
}%
\caption{\textbf{Two places to encode a watermark.}
Token-local methods alter sampled identities and verify from $y,K$ alone. Shibboleth instead
encodes which context response the checkpoint favours, so verification reruns the same $\theta$.}
\label{fig:paradigm}
\end{figure}


Nearly all language-model watermarks place the secret in token identities: the key biases a
vocabulary subset or couples randomness to each draw. Their detector therefore needs only the
text and key (Figure~\ref{fig:paradigm}a)~\citep{kirchenbauer2023watermark,
kuditipudi2024distortionfree}. This design is inexpensive, but it must coexist with the decoder.
For masked DLMs, that decoder may resolve many positions in parallel and choose their order from
confidence, so forcing a resolved prefix or steering the order can change the system being
watermarked~\citep{hong2026dgmark}.

This tension suggests a different question: can the originating model, rather than the token alone,
carry the signal? A masked DLM can answer a reproducible question about its own output. Take a
completed sentence, hide one subset of its earlier context, and record the
probability of an emitted token; then hide a second subset and record it again. Some candidates
become more likely and others less likely. The sign of that change is a \emph{model response}.
Because the verifier can recreate both hidings from the finished text, this response can carry a
keyed signal even when the token identity alone does not (Figure~\ref{fig:paradigm}b).

Two obstacles separate this observation from a valid detector. First, verification must recreate
the generation-time state even though the full continuation is now visible. Second, the method may
seek positions where both response directions are plausible, but it must finish that search before
the key reveals which direction counts as a match. Reversing this order would bias the selected
positions toward the requested answer.

Shibboleth addresses both obstacles. For each active block, it re-masks that block and its suffix,
compares pairs of hidden-prefix probes, and selects positions whose smaller directional probability
total exceeds a fixed cutoff. The key then supplies the requested direction. At a selected position,
Shibboleth draws up to $R$ candidates from the decoder's current distribution and commits the first
match; all other positions follow the reference decoder. Verification reconstructs the probes,
repeats the selection rule, and evaluates an exact binomial tail. The block partition, confidence
order, and checkpoint probabilities remain those of the reference decoder.

Our contributions are:
\begin{itemize}[leftmargin=1.2em,itemsep=1pt,topsep=2pt]
\item We introduce model-response watermarking: the signal is a checkpoint's reproducible reaction
to hidden context rather than a keyed pattern in token identities.
\item We give a replayable construction whose carrier positions are selected before the keyed
orientation is consulted, yielding finite-sample false-positive control while preserving the
reference decoding schedule (Sections~\ref{sec:setup}--\ref{sec:theory}).
\item We evaluate two masked DLMs. At $1\%$ FPR, Shibboleth reaches $0.80$ TPR on LLaDA and
$0.90$ on Dream. Verification requires 144 full-sequence model calls, and detection power
decreases under dispersed edits (Section~\ref{sec:experiments}).
\end{itemize}
